\begin{abstract}
Modeling financial time series (FTS) remains a central challenge in Quantitative Finance, since returns exhibit stylized facts such as heavy tails, volatility clustering, and the leverage effect that break the assumptions of classical models. While score-based diffusion models have reached state-of-the-art generation in other domains, their application to FTS is still limited. This thesis investigates whether diffusion models can learn the distributional properties of S&P 500 log-returns and
reproduce their core stylized facts.

The work is structured in two sequential phases. Phase 1 replicates the approach of
Kim et al., implementing VE-, VP-, and GBM-inspired SDE
variants on univariate Close log-returns to establish a validated baseline and to
expose the design ambiguities present in the original paper. Phase 2 embeds a
Conditional Score-based Diffusion (CSDI) backbone within the Elucidated Diffusion
Model (EDM) framework of Karras et al., and introduces
multivariate conditioning, learning to generate Close log-returns given the
contemporaneous Open, High, and Low observations. Models are evaluated against
GARCH-X and Merton-X parametric baselines using distributional similarity measures
($D_{KS}$, $W_1$, tail exponent $\hat{\alpha}$), probabilistic forecasting measures
(CRPS and interval coverage), and an assessment of the stylized facts.

The results show that the CSDI-based architecture can learn the bulk of the FTS
distribution under the SDE framework, but consistently underestimates higher
moments and tail heaviness, while the GBM family in its current formulation fails outright. Only the combination of EDM preconditioning and same-day Open, High, and Low conditioning
substantially narrows the gap to the empirical distribution: the excess-kurtosis
deficit shrinks from $49\%$ to $1\%$ and $W_1$ improves roughly fivefold, while
conditioning yields a CRPS skill score of about $0.74$ over the climatology
baseline. More broadly, the work generalizes the design framework to enable easier
replicability and a higher degree of control over training choices.
\end{abstract}